% =====================================================================
% A Hybrid Post-Quantum WireGuard with Cryptographic Agility:
% Instantiation with the SM Suite and Flexible KEM Switching
%
% Author: Chen Xiaochuan
% Date  : 2026-06
%
% Formatted according to the Science China Information Sciences
% (SCIS) template (SCIS2026.cls).
% =====================================================================
\documentclass{SCIS2026}

%%%%%%%%%%%%%%%%%%%%%%%%%%%%%%%%%%%%%%%%%%%%%%%%%%%%%%%
%%% Author's definitions for this manuscript
%%%%%%%%%%%%%%%%%%%%%%%%%%%%%%%%%%%%%%%%%%%%%%%%%%%%%%%

% ---------- Extra packages not provided by SCIS2026.cls --------------
% SCIS2026.cls already loads most required packages.
% xspace is needed for the convenience macros defined below.
\usepackage{xspace}

% Indent the first paragraph of every (sub)section too; by default LaTeX
% suppresses indentation on the first paragraph after a heading.
\usepackage{indentfirst}

% Figures (per-suite agility evaluation). graphicx may already be loaded by
% SCIS2026.cls; \usepackage is idempotent for options-free loads.
% Images are uploaded to the Overleaf project root, so no \graphicspath needed.
\usepackage{graphicx}

% SCIS2026.cls already loads algorithm + algorithmic;
% provide a few convenience wrappers for readability.
\newcommand{\Statex}[0]{\item[]}
\algsetup{indent=2em}

% ---------- Code listings -------------------------------------------
\usepackage{listings}
\usepackage{xcolor}
\definecolor{codebg}{rgb}{0.96,0.96,0.96}
\definecolor{codekey}{rgb}{0.0,0.0,0.55}
\definecolor{codecmt}{rgb}{0.20,0.45,0.20}
\definecolor{codestr}{rgb}{0.55,0.0,0.0}
\lstdefinelanguage{rust}{
  morekeywords={fn,let,mut,pub,struct,enum,trait,impl,for,if,else,match,
                use,mod,return,as,const,static,Self,self,type,unsafe,
                where,move,ref,while,loop,break,continue,in,crate,extern,
                async,await,dyn,box,macro_rules},
  morekeywords=[2]{u8,u16,u32,u64,usize,i8,i16,i32,i64,bool,str,String,
                   Vec,Option,Result,Some,None,Ok,Err,GenericArray,Cell},
  sensitive=true,
  morecomment=[l]{//},
  morecomment=[s]{/*}{*/},
  morestring=[b]",
  morestring=[b]',
}
\lstset{
  basicstyle=\ttfamily\scriptsize,
  backgroundcolor=\color{codebg},
  keywordstyle=\color{codekey}\bfseries,
  keywordstyle=[2]\color{codekey},
  commentstyle=\color{codecmt}\itshape,
  stringstyle=\color{codestr},
  showstringspaces=false,
  columns=fullflexible,
  breaklines=true,
  breakatwhitespace=true,
  frame=single,
  framesep=2pt,
  framerule=0.3pt,
  rulecolor=\color{black!30},
  xleftmargin=2pt,
  xrightmargin=2pt,
  belowskip=4pt,
  aboveskip=4pt,
  language=rust,
}

% ---------- Convenient macros ---------------------------------------
\newcommand{\smtwo}{\textsf{SM2}\xspace}
\newcommand{\smthree}{\textsf{SM3}\xspace}
\newcommand{\smfour}{\textsf{SM4}\xspace}
\newcommand{\xcurve}{\textsf{X25519}\xspace}
\newcommand{\blake}{\textsf{BLAKE2s}\xspace}
\newcommand{\chacha}{\textsf{ChaCha20--Poly1305}\xspace}
\newcommand{\xchacha}{\textsf{XChaCha20--Poly1305}\xspace}
\newcommand{\sgcm}{\textsf{SM4-GCM}\xspace}
\newcommand{\sctr}{\textsf{SM4-CTR}\xspace}
\newcommand{\mlkem}{\textsf{ML-KEM-512}\xspace}
\newcommand{\mceliece}{\textsf{Classic-McEliece-460896}\xspace}
\newcommand{\hybridwg}{\textsf{Hybrid-WG}\xspace}
\newcommand{\sysname}{\textsf{SM-PQ-WireGuard}\xspace}
\newcommand{\initmsg}{\textsc{InitHello}\xspace}
\newcommand{\respmsg}{\textsc{RespHello}\xspace}
\newcommand{\concat}{\mathbin{\Vert}}
\newcommand{\negl}{\mathsf{negl}}
\newcommand{\Adv}{\mathsf{Adv}}
\newcommand{\Exp}{\mathsf{Exp}}
\newcommand{\Hash}{\mathsf{H}}
\newcommand{\HMAC}{\mathsf{HMAC}}
\newcommand{\HKDF}{\mathsf{HKDF}}
\newcommand{\KDF}{\mathsf{KDF}}
\newcommand{\AEAD}{\mathsf{AEAD}}
\newcommand{\KEMs}{\textsf{KEM}_s}
\newcommand{\KEMe}{\textsf{KEM}_e}
\newcommand{\ECDH}{\mathsf{ECDH}}
\newcommand{\Encap}{\mathsf{Encap}}
\newcommand{\Decap}{\mathsf{Decap}}
\newcommand{\KeyGen}{\mathsf{KeyGen}}
\newcommand{\pkenc}{\mathsf{pkenc}}
% Provide \qed if not already defined
\providecommand{\qed}{\hfill$\square$}

%%%%%%%%%%%%%%%%%%%%%%%%%%%%%%%%%%%%%%%%%%%%%%%%%%%%%%%
%%% Begin
%%%%%%%%%%%%%%%%%%%%%%%%%%%%%%%%%%%%%%%%%%%%%%%%%%%%%%%
\begin{document}

%%%%%%%%%%%%%%%%%%%%%%%%%%%%%%%%%%%%%%%%%%%%%%%%%%%%%%%
%%% Article information
%%%%%%%%%%%%%%%%%%%%%%%%%%%%%%%%%%%%%%%%%%%%%%%%%%%%%%%
\ArticleType{RESEARCH PAPER}
\Year{2026}
\Month{June}
\Vol{69}
\No{1}
\DOI{}
\ArtNo{}
\ReceiveDate{}
\ReviseDate{}
\AcceptDate{}
\OnlineDate{}
\AuthorMark{}
\AuthorCitation{}

%%% title
\title{A Hybrid Post-Quantum WireGuard with Cryptographic Agility:
       Instantiation with the SM Suite and Flexible KEM Switching}
       {A Hybrid Post-Quantum WireGuard with Cryptographic Agility}

%%% authors
%%% Corresponding author: include email in double braces {{email@xxx.com}}
%%% General author: leave third argument empty {}
\author[1]{Xiaochuan CHEN}{}
\author[1]{Yifan ZHENG}{}
\author[1]{Taolong SU}{}
\author[2]{Feng ZHANG}{}
\author[1]{Zhe LIU}{{zhe.liu@zju.edu.cn}}

%%% Address (please fill in actual affiliations)
\address[1]{School of Software Technology, Zhejiang University, Ningbo 315100, China}
\address[2]{Hangzhou Post-quantum Cryptography Technology Co., Ltd}

%%% Abstract
\abstract{The advent of quantum computers poses a significant threat to
classical VPN protocols such as WireGuard, whose security relies on the
hardness of the elliptic-curve discrete logarithm problem. Moreover, any
single post-quantum key encapsulation mechanism may be subject to future
cryptanalytic breakthroughs or prove suboptimal in specific deployment
scenarios, making cryptographic agility a critical requirement for
deployable solutions. The GM/T family of cryptographic standards
(SM2, SM3, SM4) has been established as a general-purpose national
cryptographic suite and is widely encouraged for use in information
systems. To address these challenges, we present SM-PQ-WireGuard, the
first hybrid post-quantum WireGuard handshake that seamlessly unifies
GM/T compliance with cryptographic agility. The protocol instantiates
the classical layer with the SM2 elliptic curve, SM3 hash, and SM4
symmetric cipher, and pairs a fixed Classic-McEliece-460896 static-KEM
identity anchor with an agile ephemeral-KEM slot supporting a range of
post-quantum options including ML-KEM, Kyber, HQC, and FrodoKEM. An
algorithm-agnostic abstraction layer allows flexible switching of the
ephemeral KEM without modifying the core handshake logic, enabling
operators to tailor the protocol to diverse deployment constraints while
preserving full GM/T compliance. We provide a rigorous
security analysis in the random oracle model, reducing the security of
SM-PQ-WireGuard to the gap Diffie--Hellman assumption on SM2 and the
IND-CCA security of the underlying KEMs. This computational proof is
complemented by mechanised symbolic verification---ProVerif for trace and
equivalence properties, independently cross-checked by Tamarin---which
confirms that all supported KEM configurations preserve the same
hybrid-security properties. Protocol-layer adaptations including MAC1/MAC2
reconstruction via HMAC-SM3 truncation and per-site nonce partitioning
for SM4-GCM are explicitly derived to address the semantic gaps
introduced by the SM primitive suite. To the best of our knowledge, this
is the first work to present a complete GM/T-compliant post-quantum
WireGuard handshake with cryptographic agility. Our implementation
demonstrates that the GM/T substitution and the agility framework incur
minimal overhead, validating the practicality of the design.}

%%% Keywords
\keywords{cryptographic agility, hybrid post-quantum cryptography, WireGuard, SM suite}

\maketitle

% =====================================================================
\section{Introduction}
\label{sec:intro}

WireGuard~\cite{1} has become one of the most
widely deployed VPN protocols, supplanting legacy IPsec installations
in both cloud and enterprise environments.  Its appeal stems from a
small attack surface, a single round-trip handshake, and a key exchange
based on the Noise framework~\cite{2}, itself rooted in
Diffie--Hellman key exchange~\cite{3}.  The current
handshake instantiates Noise \textsf{IKpsk2} with
X25519~\cite{4},
ChaCha20--Poly1305~\cite{5}, and
BLAKE2s~\cite{6}, deriving a session key that protects
every IP packet between peers.  However, the security of X25519 rests
on the Elliptic-Curve Discrete Logarithm Problem~\cite{7,8};
a sufficiently large quantum computer running Shor's
algorithm~\cite{9} would break it retroactively.  This
``harvest-now, decrypt-later'' threat~\cite{10} has
motivated hybrid key
exchange~\cite{11,12,13},
where the session key combines a classical Diffie--Hellman secret with
one or more post-quantum KEM secrets, ensuring that breaking one world
does not compromise the other~\cite{14}.

Most existing hybrid proposals hard-code a specific KEM
selection---typically ML-KEM-512~\cite{15} for ephemeral
exchange and Classic-McEliece-460896~\cite{16} for
static protection.  This one-size-fits-all
approach overlooks a fundamental reality: \emph{no single KEM is
optimal for every environment}.  Bandwidth-constrained networks favour
small ciphertexts; storage-constrained devices favour small public
keys; latency-sensitive applications favour fast encapsulation; and
future cryptanalytic advances may render any single choice obsolete.
The GM/T family of cryptographic standards---SM2~\cite{17,18},
SM3~\cite{19}, and SM4~\cite{20}---is established as a
general-purpose national cryptographic suite and is widely encouraged
for use in Chinese information systems.  A deployable post-quantum VPN
for the Chinese market must therefore satisfy three simultaneous
requirements: quantum resistance, GM/T standards compliance, and
cryptographic agility.

To address the aforementioned issues, we present SM-PQ-WireGuard, a
hybrid post-quantum WireGuard handshake that unifies GM/T standards
compliance with cryptographic agility.  The protocol instantiates the
classical layer with SM2 ECDH, SM3 hashing, and SM4-GCM/SM4-CTR
symmetric encryption, and combines two KEM slots: a static slot
\emph{fixed} to Classic-McEliece-460896 as the long-term identity
anchor, and an agile ephemeral slot switchable among
ML-KEM-512/768/1024, Kyber-512, HQC-128, and Frodo-640-AES.  An
algorithm-agnostic abstraction layer decouples core handshake logic
from the concrete ephemeral KEM; a one-byte suite identifier hashed into the handshake
transcript ensures distinct session keys across algorithm choices,
preventing cross-algorithm confusion.  This design allows operators to
select KEMs based on deployment constraints without protocol-level
changes or security-proof invalidation.

Beyond the agility framework, we address two protocol-layer mismatches
that arise from the SM primitive suite: (i)~SM3 lacks the native
keyed-hash and variable-output modes that upstream WireGuard relied on
for MAC1/MAC2, which we reconstruct via HMAC-SM3 truncation; and
(ii)~SM4-GCM's 96-bit nonce demands a per-site strategy across four
protocol locations where upstream used extended nonces.  We explicitly
re-derive the Noise transcript for the GM/T setting to resolve these
gaps, ensuring correctness and security.  We provide a rigorous
IND-CCA security reduction in the random oracle model, complemented by
mechanised symbolic verification (ProVerif, cross-checked by Tamarin)
confirming that every ephemeral-KEM configuration preserves
the same hybrid-security guarantees.  Performance evaluation on a
commodity x86\_64 host demonstrates that the GM/T substitution and
agility framework incur minimal overhead, validating the practicality
of the design.

We make four contributions.
\begin{itemize}
\item We design SM-PQ-WireGuard, a Noise \textsf{IKpsk2}-compliant
protocol that unifies SM2/SM3/SM4 with a fixed Classic-McEliece static
identity anchor and an agile ephemeral-KEM slot, supporting flexible
switching across a range of options (ML-KEM, Kyber, HQC, FrodoKEM)
without protocol-level changes.

\item We explicitly re-derive the Noise transcript for the GM/T
setting, resolving two semantic gaps introduced by the primitive
substitution: MAC1/MAC2 reconstruction via HMAC-SM3 truncation, and
per-site nonce partitioning for SM4-GCM across four distinct protocol
locations.

\item We provide an IND-CCA security reduction under standard
assumptions on SM2 and SM3, complemented by mechanised symbolic
verification (ProVerif, independently cross-checked by Tamarin) that
every supported ephemeral KEM maintains identical hybrid-security
guarantees.

\item We implement SM-PQ-WireGuard with a feature-gated benchmark
harness.  The implementation validates that the GM/T substitution and
agility framework incur minimal overhead, demonstrating the
practicality of a compliant and flexible post-quantum VPN design.
\end{itemize}

\paragraph{Roadmap.}
Section~\ref{sec:bg} recalls WireGuard, the Noise framework, the
GM/T suite, and the post-quantum KEM options.
Section~\ref{sec:design} presents SM-PQ-WireGuard in detail, including
the agility framework and protocol-layer adaptations.
Section~\ref{sec:security} contains the dual-track security analysis
(computational and symbolic).  Section~\ref{sec:perf} reports the
performance evaluation.  Section~\ref{sec:disc} discusses deployment
considerations, limitations and future work, and concludes the paper.
Proofs are in the appendix.


% =====================================================================
\section{Background}
\label{sec:bg}

\subsection{Notation}
\label{sec:notation}

Throughout the paper, $\lambda$ denotes the security parameter; we
target $\lambda = 128$ bits of classical security.  A function $f$
is \emph{negligible}, written $f(\lambda) = \negl(\lambda)$, if for
every polynomial $p$ there exists $\lambda_0$ such that
$f(\lambda) < 1/p(\lambda)$ for all $\lambda \geq \lambda_0$, following
standard conventions~\cite{21}.  For
bit strings $x$ and $y$, $|x|$ denotes the bit length of $x$ and
$x \concat y$ denotes concatenation.  $x \xleftarrow{\$} S$ denotes
uniform random sampling from set $S$.  We write $\KDF_1, \KDF_2,
\KDF_3$ for the Noise-style key derivation variants producing one,
two, or three $256$-bit outputs from a chaining key and input key
material.

\subsection{WireGuard and the Noise Framework}
\label{sec:wireguard}

WireGuard's handshake is a syntactic specialisation of the Noise
\textsf{IKpsk2} pattern~\cite{2}:
\begin{align*}
\text{IKpsk2} :=\;& \leftarrow s; \\
              & \to e, es, s, ss; \\
              & \leftarrow e, ee, se,\ \mathsf{psk}.
\end{align*}
The initiator knows the responder's static public key in advance.
In the first flight, the initiator sends an ephemeral public key,
performs an ephemeral-static DH, sends its own static public key
(encrypted under a key derived from the transcript), and performs a
static-static DH.  The responder replies with its ephemeral public
key, performs ephemeral-ephemeral and static-ephemeral DHs, and
mixes in a pre-shared key.  The final chaining key is split into
send and receive keys for the data plane.

Stock WireGuard instantiates this pattern with X25519~\cite{4}
for ECDH, ChaCha20--Poly1305~\cite{5} for
AEAD~\cite{22,23}, and
BLAKE2s~\cite{6} for hashing and key derivation.
The session key is derived from the concatenation of a classical DH
secret and, in hybrid variants, one or more post-quantum KEM secrets,
such that breaking only the classical component does not compromise
the session.  We adopt this hybrid structure with two KEM slots: one
static (for long-term identity protection) and one ephemeral (for
per-session freshness).

\subsection{The GM/T Cryptographic Suite}
\label{sec:gmt}

The GM/T suite comprises three standardised algorithms widely adopted
in Chinese information systems.

\textbf{SM2~\cite{17,18}} is an elliptic-curve
cryptosystem defined over a $256$-bit prime field with a short
Weierstrass equation.  It has prime order and cofactor $1$.  A
private key is a scalar $d \in [1, n-1]$; the public key is the
affine point $P = d\cdot G$.  SEC1 compressed encoding~\cite{24} represents a
point as $y_{\text{parity}} \concat X$, where $y_{\text{parity}} \in
\{\texttt{0x02},\texttt{0x03}\}$ and $X$ is the $32$-byte
$x$-coordinate, yielding a $33$-byte public key.  The shared secret
in SM2 ECDH is the $32$-byte $x$-coordinate of $d_A \cdot d_B \cdot
G$.  Security rests on the gap Diffie--Hellman assumption on the SM2
curve.

\textbf{SM3~\cite{19}} is a $256$-bit hash function following
the Merkle--Damg\aa rd construction~\cite{25,26}
with a $64$-byte input block.  Unlike BLAKE2s, SM3 provides no native
keyed-hash mode and no variable-length output.  This motivates the
MAC1/MAC2 reconstruction we present in Section~\ref{sec:macreconstruct};
its security margin remains intact under published
cryptanalysis.

\textbf{SM4~\cite{20}} is a $128$-bit block cipher with a
$128$-bit key, employing a $32$-round unbalanced Feistel network whose
security has been examined by dedicated cryptanalysis.
We use SM4 in GCM mode~\cite{27,28} (with a $96$-bit
nonce) for AEAD and in CTR mode for stream encryption.  The fixed $96$-bit nonce---half the
size of XChaCha20--Poly1305's $192$-bit nonce---requires per-site
nonce partitioning, which we address in
Section~\ref{sec:noncedomain}.

\subsection{Post-Quantum KEM Options}
\label{sec:kemoptions}

To support cryptographic agility, SM-PQ-WireGuard accommodates a
range of post-quantum KEMs in its \emph{ephemeral} slot, while the
static slot is fixed to Classic-McEliece-460896~\cite{16} as the
long-term identity anchor.  Classic-McEliece is well suited to this
role: its conservative, long-studied code-based security matches the
longevity of an identity key, and although its public key is large
($524{,}160$\,B, $\approx512$\,KiB) it is generated once and pinned out
of band rather than transmitted each handshake, while its
ciphertext---which \emph{is} sent in \initmsg---is among the smallest of
all candidates ($156$\,B).
Table~\ref{tab:kems} summarises the agile ephemeral
candidates---ML-KEM~\cite{15}, Kyber~\cite{29}, HQC, and
FrodoKEM---and their key characteristics.  These KEMs are selected to represent diverse
security assumptions (lattice-based (ML-KEM, Kyber), code-based (HQC), and unstructured-LWE (FrodoKEM) designs) and performance trade-offs, allowing
operators to choose the most suitable ephemeral configuration for their
deployment constraints.

\begin{table}[htbp]
\centering
\caption{Supported agile ephemeral-KEM candidates and their key
characteristics (the static slot is fixed to Classic-McEliece-460896).
$|\textsf{pk}|$ and $|\textsf{ct}|$ denote public
key and ciphertext sizes respectively.}
\label{tab:kems}
\footnotesize
\tabcolsep 5pt
\begin{tabular}{llrrl}
\toprule
\textbf{KEM} & \textbf{Category} & $|\textsf{pk}|$ & $|\textsf{ct}|$ & \textbf{Security} \\
\midrule
ML-KEM-512    & Lattice            & $800$\,B  & $768$\,B  & AES-128 \\
ML-KEM-768    & Lattice            & $1184$\,B & $1088$\,B & AES-192 \\
ML-KEM-1024   & Lattice            & $1568$\,B & $1568$\,B & AES-256 \\
Kyber-512     & Lattice            & $800$\,B  & $768$\,B  & AES-128 \\
HQC-128       & Code-based         & $2249$\,B & $4433$\,B & AES-128 \\
Frodo-640-AES & Unstructured lattice (LWE) & $9616$\,B & $9720$\,B & AES-128 \\
\bottomrule
\end{tabular}
\end{table}


% =====================================================================
\section{Design of SM-PQ-WireGuard}
\label{sec:design}

We present the design of SM-PQ-WireGuard, a hybrid post-quantum
WireGuard handshake that unifies GM/T compliance with cryptographic
agility.  The protocol follows the Noise \textsf{IKpsk2} pattern,
instantiates the classical layer with the SM2/SM3/SM4 suite, and
combines a fixed Classic-McEliece static KEM anchor with an agile
ephemeral KEM slot supporting flexible switching without protocol-level
changes.  This section describes the design goals, the protocol
overview, the handshake in detail, the protocol-layer adaptations
required by the SM suite, and the agility framework that enables KEM
switching.

\subsection{Design Goals}
\label{sec:goals}

The design of SM-PQ-WireGuard is guided by four goals.

\begin{description}
\item[G1: Semantic equivalence.] The protocol must be a syntactic
specialisation of Noise \textsf{IKpsk2}, differing only in primitive
choices and wire encodings.

\item[G2: Transcript unambiguity.] Sessions must never be replayable
against other WireGuard variants or intermediate constructions; this
is enforced by a unique protocol identifier that acts as a domain
separator.

\item[G3: Cryptographic agility.] Operators must be able to select
and switch the ephemeral KEM without modifying the core handshake logic
or invalidating security proofs; the static slot is a fixed
Classic-McEliece identity anchor.

\item[G4: Wire-format minimality.] The wire format must change
minimally---the only difference from a non-SM hybrid construction is
the $33$-byte SEC1-compressed SM2 ephemeral public key replacing the
$32$-byte X25519 $u$-coordinate.  With the default ML-KEM-512 (or
Kyber-512) ephemeral suite, both handshake messages remain within the
IPv6 MTU of $1280$ bytes; the higher-security and conservative-family
suites (ML-KEM-768/1024, HQC-128, Frodo-640-AES) trade larger
ciphertexts for stronger or more diverse assumptions and exceed the
MTU, so they rely on IP fragmentation (\S\ref{sec:perf}).
\end{description}

\subsection{Protocol Overview}
\label{sec:overview}

SM-PQ-WireGuard instantiates the Noise \textsf{IKpsk2} pattern with
the GM/T cryptographic suite at the classical layer: SM2 for ECDH,
SM3 for hashing and key derivation, and SM4-GCM/SM4-CTR for
symmetric encryption.  The handshake incorporates two post-quantum
KEM slots: a static slot for long-term identity protection and an
ephemeral slot for per-session freshness.  An algorithm-agnostic
abstraction layer decouples the core handshake logic from concrete
KEM implementations, and a one-byte suite identifier
(\texttt{f\_suite}) is hashed into the handshake transcript
to ensure domain separation across different KEM choices.  The
handshake completes in two messages---\initmsg\ and
\respmsg---and produces a $256$-bit session key, truncated to
$128$ bits for SM4 at the data plane.

Figure~\ref{fig:flow} gives a high-level view of this exchange
before the formal treatment in Section~\ref{sec:handshake}.  Both
parties first generate fresh ephemeral material: the initiator
draws an SM2 ephemeral key pair \emph{and} an ephemeral
post-quantum KEM key pair, while the responder draws an SM2
ephemeral key pair, its KEM contribution being realised as an
encapsulation against the public key it receives.  The initiator's
\initmsg\ carries its SM2 and PQ ephemeral public keys together
with encrypted static identity material; from it the responder
parses every public value it needs and computes the SM2 DH and the
PQ KEM shared secrets.  The responder's \respmsg\ returns the PQ
ciphertexts and its own encrypted identity, so that the initiator
can decapsulate the same KEM secrets and reconstruct the identical
SM2 DH value.  Each side then feeds the SM2 DH secret, the PQ KEM
shared secrets, and the pre-shared key $Q$ into the SM3-based key
schedule, yielding the directional session keys that protect the
SM4-GCM data plane.  Because the post-quantum secrets enter the
same chaining key as the classical DH output, the session stays
secure as long as \emph{either} the SM2 or the KEM contribution
remains uncompromised---the hybrid guarantee made precise in
Section~\ref{sec:security}.

\begin{figure*}[htbp]
\centering
\includegraphics[width=\linewidth]{flowchart.png}
\caption{End-to-end message flow of the SM-PQ-WireGuard handshake.
The initiator and responder exchange two messages---\initmsg\ and
\respmsg---each combining SM2 (classical) ephemeral key agreement
with post-quantum KEM encapsulation.  Both parties derive the
session keys from the SM2 DH secret, the PQ KEM shared secrets, and
the pre-shared key (PSK) through the SM3 key schedule; the
resulting keys protect the SM4-GCM data plane.  The detailed
message construction is formalised in Algorithms~\ref{alg:init}
and~\ref{alg:resp}.}
\label{fig:flow}
\end{figure*}

\subsection{Handshake Protocol}
\label{sec:handshake}

The key derivation follows the Noise KDF structure.  For a chaining
key $ck \in \{0,1\}^{256}$ and input key material $m$, the
derivation produces $n\in\{1,2,3\}$ consecutive $256$-bit outputs:
\[
\KDF_n(ck, m) = (t_1, \dots, t_n),
\]
where $t_0 = \textsf{HMAC-SM3}_{ck}(m)$ and
$t_i = \textsf{HMAC-SM3}_{t_0}(t_{i-1} \concat \langle i \rangle)$
for $i=1,\dots,n$.  All HMAC calls are instantiated with SM3.  Every
AEAD key is obtained by taking the leading $128$ bits of the
corresponding $256$-bit KDF output.

\subsubsection{Initiator Flow}

Algorithm~\ref{alg:init} formalises the construction of the
\initmsg.  The initiator begins by hashing the responder's static
public keys into the transcript, generates a fresh SM2 ephemeral key
pair and a fresh ephemeral KEM key pair, and mixes both public keys
into the chaining key.  It then encapsulates against the responder's
static KEM public key to obtain $\textsf{shk}_1$, computes the SM2
ephemeral-static DH to obtain $z_1$, and wraps the static KEM
ciphertext with a stream cipher keyed from $z_1$---this sanitises
the transcript by preventing an adversary from planting chosen
encapsulation metadata.  The initiator then derives a symmetric key
$k_S$ from $z_1 \concat \textsf{shk}_1$, encrypts its identity hash,
performs the static-static SM2 DH to obtain $z_2$, hashes the two
static KEM public keys together, and mixes $z_2$, the PQ identity
hash, and the pre-shared key $Q$ into the chaining key.  A final
AEAD block carries a strictly increasing replay counter.

\begin{algorithm}[htbp]
\caption{\initmsg\ construction (initiator $i$).}
\label{alg:init}
\begin{algorithmic}[1]
\REQUIRE own static key material $(S^i, S^i_q, \mathit{sk}_{S^i}, \mathit{sk}_{S^i_q})$;
peer's static public keys $(S^r, S^r_q)$; pre-shared $Q$.
\STATE $ck \gets ck_0;\ hs \gets hs_0$
\STATE $hs \gets \textsf{SM3}(hs \concat \textsf{SM3}(S^r \concat S^r_q))$
\STATE $hs \gets \textsf{SM3}(hs \concat \texttt{f\_suite})$
\STATE $(E, e) \gets \KeyGen(\textsf{SM2})$
\STATE $(E_q, e_q) \gets \KeyGen(\KEMe)$
\STATE $ck \gets \KDF_1(ck, \pkenc(E) \concat \pkenc(E_q))$
\STATE $(\textit{ct}_1, \textsf{shk}_1) \gets \Encap(\KEMs; S^r_q)$
\STATE $z_1 \gets \ECDH(e, S^r)$
\STATE $k_0 \gets \KDF_1(0^{256}, z_1)$
\STATE $\textit{cct}_1 \gets \mathsf{Stream}(\textit{ct}_1)$
\STATE $hs \gets \textsf{SM3}(hs \concat \pkenc(E) \concat \pkenc(E_q))$
\STATE $(ck, k_S) \gets \KDF_2(ck, z_1 \concat \textsf{shk}_1)$
\STATE $c_S \gets \AEAD.\mathsf{Enc}
        \bigl(\textsf{SM3}(\pkenc(S^i) \concat \pkenc(S^i_q));\ hs\bigr)$
\STATE $hs \gets \textsf{SM3}(hs \concat c_S)$
\STATE $z_2 \gets \ECDH(\mathit{sk}_{S^i}, S^r)$
\STATE $h_{\text{pq}} \gets \textsf{SM3}(S^i_q \concat S^r_q)$
\STATE $(ck, k_T) \gets \KDF_2(ck, z_2 \concat h_{\text{pq}} \concat Q)$
\STATE $c_T \gets \AEAD.\mathsf{Enc}(\mathsf{ctr}^{16};\ hs)$
\STATE $hs \gets \textsf{SM3}(hs \concat c_T)$
\STATE \RETURN $\initmsg = \langle \pkenc(E), \pkenc(E_q),
       \textit{cct}_1, c_S, c_T\rangle$
\end{algorithmic}
\end{algorithm}

\subsubsection{Responder Flow}

Algorithm~\ref{alg:resp} formalises the responder's processing of
\initmsg\ and construction of \respmsg.  The responder generates a
fresh SM2 ephemeral key pair, encapsulates against the initiator's
ephemeral KEM public key (obtaining $\textsf{shk}_2$) and against
the initiator's static KEM public key (obtaining $\textsf{shk}_3$),
and performs two SM2 ECDH operations: ephemeral-static with the
initiator's static key ($z_3$) and ephemeral-ephemeral with the
initiator's ephemeral key ($z_4$).  The responder wraps the static
KEM ciphertext using a stream cipher keyed from $z_3$, then mixes
$z_4 \concat \textsf{shk}_2$ and $z_3 \concat \textsf{shk}_3$ into
the chaining key through successive KDF calls.  The PSK $Q$ is mixed
in via $\KDF_3$, producing a transcript commitment $\tau$ that is
hashed into $hs$ and a symmetric key $k_E$ that seals an empty AEAD
block as key confirmation.  The final $\KDF_2$ splits the chaining
key into the two data-plane keys, each truncated to $128$ bits for
use with SM4.

\begin{algorithm}[htbp]
\caption{\respmsg\ construction (responder $r$).}
\label{alg:resp}
\begin{algorithmic}[1]
\REQUIRE own static keys $(S^r, S^r_q, \mathit{sk}_{S^r}, \mathit{sk}_{S^r_q})$;
peer's ephemeral public key $E$ from \initmsg;
peer's static public keys $(S^i, S^i_q)$; pre-shared $Q$.
\STATE $(E', e') \gets \KeyGen(\textsf{SM2})$
\STATE $(\textit{ct}_2, \textsf{shk}_2) \gets \Encap(\KEMe; E_q)$
\STATE $(\textit{ct}_3, \textsf{shk}_3) \gets \Encap(\KEMs; S^i_q)$
\STATE $z_3 \gets \ECDH(e', S^i)$
\STATE $z_4 \gets \ECDH(e', E)$
\STATE $k_0 \gets \KDF_1(0^{256}, z_3)$
\STATE $\textit{cct}_3 \gets \mathsf{Stream}(\textit{ct}_3)$
\STATE $ck \gets \KDF_1(ck, \pkenc(E') \concat \textit{ct}_2)$
\STATE $hs \gets \textsf{SM3}(hs \concat \pkenc(E') \concat \textit{ct}_2)$
\STATE $ck \gets \KDF_1(ck, z_4 \concat \textsf{shk}_2)$
\STATE $ck \gets \KDF_1(ck, z_3 \concat \textsf{shk}_3)$
\STATE $(ck, \tau, k_E) \gets \KDF_3(ck, Q)$
\STATE $hs \gets \textsf{SM3}(hs \concat \tau)$
\STATE $c_E \gets \AEAD.\mathsf{Enc}(\varepsilon;\ hs)$
\STATE $(k_{r\to i}, k_{i\to r}) \gets \KDF_2(ck, \varepsilon)$
\STATE \RETURN $\respmsg = \langle \pkenc(E'), \textit{ct}_2,
       \textit{cct}_3, c_E\rangle$
\end{algorithmic}
\end{algorithm}

\subsection{Protocol Adaptations for the SM Suite}
\label{sec:adaptations}

The SM primitive suite introduces two protocol-layer mismatches that
a naive cipher swap would miss.  We resolve both through explicit
transcript re-derivation.

\subsubsection{MAC1/MAC2 Reconstruction via HMAC-SM3 Truncation}
\label{sec:macreconstruct}

Every WireGuard handshake message carries two $16$-byte MAC fields:
MAC1 for cookie-resolved sender authentication and MAC2 for
stateless DoS protection.  Stock WireGuard derives both using
BLAKE2s in its native keyed-hash, variable-output mode: a single
library call with a key and a desired output length of $16$ bytes
directly produces the tag.

SM3 provides no native keyed mode and no variable-length output.
We therefore reconstruct both MACs using the HMAC construction
(RFC~2104)~\cite{30} instantiated with SM3, then truncate
the $256$-bit HMAC
output to $128$ bits:
\[
\textsf{MAC}(k, m) = \textsf{HMAC-SM3}(k, m)[0\ldots 15].
\]
This replaces a single native BLAKE2s call with two HMAC compression
invocations plus an explicit truncation step.  Security follows from
the dual-PRF property of HMAC-SM3: under the random-oracle model,
the $128$-bit prefix of the $256$-bit HMAC output is
indistinguishable from a uniformly random $128$-bit string.  The
unforgeability of MAC1/MAC2 at the $128$-bit security level is
therefore preserved.

\subsubsection{Nonce-Domain Partitioning for SM4-GCM}
\label{sec:noncedomain}

SM-PQ-WireGuard employs symmetric encryption at four distinct
protocol sites, each with different nonce-management requirements.
Stock WireGuard handles these sites with three different algorithms:
ChaCha20--Poly1305 for the handshake AEAD,
XChaCha20--Poly1305 (with a $192$-bit random nonce) for cookie-reply
AEAD, bare ChaCha20 for KEM-ciphertext wrapping, and a separate
ChaCha20--Poly1305 instance for the data plane.  SM4-GCM accepts
only a $96$-bit nonce---half the size of XChaCha20--Poly1305's
$192$-bit nonce---and there is no extended-nonce analogue for SM4 in
production libraries.  We therefore assign an independent nonce
strategy to each site, summarised in Table~\ref{tab:nonces}.

\begin{table}[htbp]
\centering
\caption{Nonce strategies for the four SM4-GCM/SM4-CTR protocol sites.}
\label{tab:nonces}
\footnotesize
\tabcolsep 8pt
\begin{tabular}{lll}
\toprule
\textbf{Site} & \textbf{Algorithm} & \textbf{Nonce/IV Strategy} \\
\midrule
Handshake AEAD ($c_S, c_T, c_E$) & SM4-GCM & Fixed zero nonce (key is per-handshake fresh) \\
KEM ciphertext wrapping ($\textit{cct}_1, \textit{cct}_3$) & SM4-CTR & Fixed zero IV (key is per-handshake fresh) \\
Cookie reply AEAD & SM4-GCM & Leading $12$ bytes of $24$-byte random nonce \\
Data plane & SM4-GCM & $4$-byte zero prefix + $8$-byte monotonic counter \\
\bottomrule
\end{tabular}
\end{table}

The handshake AEAD and KEM wrapping use fixed zero nonces because
every key is freshly derived and used for exactly one message per
handshake.  The cookie reply AEAD uses the leading $12$ bytes of a
$24$-byte random nonce (preserving the upstream wire format); the
cookie key rotates every $120$ seconds and is rate-limited, bounding
the nonce-collision probability to approximately $2^{-53}$.  The
data plane uses an $8$-byte monotonic packet counter, ensuring
nonces never repeat under a single key.  This per-site partitioning
enables SM4-GCM to safely replace the three different
ChaCha20-based constructions used upstream.

\subsection{The Agility Framework}
\label{sec:agility}

To support flexible KEM switching without protocol-level changes,
SM-PQ-WireGuard introduces an algorithm-agnostic abstraction layer
with three components.

\paragraph{Abstraction layer.}
The core handshake logic interacts with KEMs solely through a trait
defining three operations: $\KeyGen \to (\textsf{pk}, \textsf{sk})$,
$\Encap(\textsf{pk}) \to (\textsf{ct}, k)$, and
$\Decap(\textsf{sk}, \textsf{ct}) \to k$.  Classic-McEliece-460896 (the
static anchor) and every ephemeral candidate---ML-KEM-512, ML-KEM-768,
ML-KEM-1024, Kyber-512, HQC-128, and Frodo-640-AES---expose this uniform
interface.  The handshake state machine (Algorithms~\ref{alg:init}
and~\ref{alg:resp}) is parameterised by the fixed static KEM
(Classic-McEliece) and a deployment-selected ephemeral KEM, without
altering the state machine logic.

\paragraph{Domain separation via the suite identifier.}
Each ephemeral-KEM suite is assigned a stable one-byte on-wire
identifier \texttt{f\_suite} (the static slot is fixed to the
Classic~McEliece identity anchor and carries no selector).  The
initiator stamps \texttt{f\_suite} into both handshake messages, and
the byte is bound into the protocol in two ways.  First, it is hashed
into the transcript hash at the start of the handshake,
\[
hs \gets \textsf{SM3}(hs \concat \texttt{f\_suite}),
\]
so a tampered selector makes the two peers' transcripts diverge
(downgrade resistance).  Second, it is covered by the PSK-keyed MAC1
(Section~\ref{sec:macreconstruct}).  Because the identifier enters the
transcript that seeds every subsequent KDF call, different KEM choices
yield distinct transcripts and session keys, preventing cross-algorithm
replay or confusion.  As the selector is a public value it does not
affect secrecy; the symbolic model (Section~\ref{sec:security}) treats
the KEM abstractly, so all supported suites inherit the same
hybrid-security guarantees.

\paragraph{Deployment and configuration.}
Operators select the ephemeral KEM via configuration at deployment
time (the static slot is fixed to Classic-McEliece).  The protocol
binary is compiled once with support for all ephemeral KEMs; the
abstraction layer dispatches the correct implementation based on the
configuration.  No code changes or recompilation are required to switch
between supported ephemeral KEMs.
This design allows seamless adaptation to different deployment
constraints---bandwidth-sensitive deployments can choose a KEM with
small ciphertexts, while latency-sensitive deployments can choose a
KEM with fast encapsulation---without re-engineering the protocol.

\subsection{Key-Length Adaptation}
\label{sec:keylen}

SM4-GCM consumes a $128$-bit key, half the length of
ChaCha20--Poly1305's $256$-bit key.  The KDF, however, continues to
operate at $256$ bits to preserve the entropy contributed by each
KEM secret and ECDH shared secret.  Truncation is applied uniformly
at the point where key material leaves the KDF and enters the AEAD
call: each $256$-bit output is truncated to its leading $128$ bits.
This does not reduce security, as the prefix of a uniformly random
$256$-bit string is itself uniformly distributed.
% =====================================================================
\section{Security Analysis}
\label{sec:security}

\subsection{Threat Model and Assumptions}
\label{sec:threat}

We analyse SM-PQ-WireGuard in the standard Dolev--Yao network model:
the adversary controls all communication channels and may intercept,
modify, inject, or replay messages.  For the computational reduction,
we additionally grant the adversary full control over session
scheduling and key-reveal queries, following the standard
real-or-random (RoR) security experiment for authenticated key
exchange~\cite{31,32}.  The adversary may compromise any
subset of long-term keys (static SM2 keys and static KEM keys) and
ephemeral keys (ephemeral SM2 keys and ephemeral KEM keys); the target
session is considered fresh if its session key has not been revealed
and at least one of the three security legs remains intact---the
classical SM2 leg (neither peer's static SM2 key is compromised and the
ephemeral SM2 keys of the session are unrevealed), the ephemeral
post-quantum leg (the ephemeral KEM key pair of the session is
uncompromised), or the static post-quantum leg (neither peer's static
KEM key is compromised).

We rely on three standard hardness assumptions: (i)~the gap
Diffie--Hellman (GDH) assumption on the SM2
curve~\cite{17,18}; (ii)~the IND-CCA security of each
supported KEM (ML-KEM-512/768/1024, Kyber-512, HQC-128, and
Frodo-640-AES); and (iii)~the random-oracle~\cite{33} / dual-PRF model for
HMAC-SM3~\cite{34}.  These assumptions are standard for
their respective primitives and are inherited from the underlying
cryptographic standards.

\subsection{Computational Security}
\label{sec:formal}

We establish the security of SM-PQ-WireGuard through a sequence of
game hops, following the standard hybrid key-exchange
reduction~\cite{32}.  The core argument is that the
session key remains pseudorandom as long as at least one of the three
underlying hardness assumptions holds.

\begin{theorem}[Hybrid security of SM-PQ-WireGuard]\label{thm:hybrid}
Let $\Pi$ be the protocol defined by Algorithms~\ref{alg:init}
and~\ref{alg:resp}, with $q_s$ the number of session queries.  Under
the GDH assumption on SM2, the IND-CCA security of the underlying
KEMs, and the random-oracle / dual-PRF assumption on HMAC-SM3, the
distinguishing advantage of any PPT adversary $\mathcal A$ is bounded by
\begin{align*}
\Adv^{\textsc{ror}}_{\Pi}(\mathcal A) \leq\
  q_s\cdot \min\bigl\{
  &\Adv^{\mathrm{GDH}}_{\textsf{SM2}}(\mathcal A_1),\
   \Adv^{\mathrm{IND\text{-}CCA}}_{\KEMe}(\mathcal A_2),\ \\
  &\Adv^{\mathrm{IND\text{-}CCA}}_{\KEMs}(\mathcal A_3)
  \bigr\}
   + \negl(\lambda),
\end{align*}
for adversaries $\mathcal A_1, \mathcal A_2, \mathcal A_3$ whose
running time is essentially that of $\mathcal A$.  The minimum over the
three terms is the precise statement of hybrid security: the advantage
stays negligible as long as the classical leg \emph{or} either
post-quantum leg is unbroken.
\end{theorem}

\begin{proof}[Proof sketch]
We model SM3 (via HMAC-SM3) as a random oracle---a standard
assumption for Noise-derived protocols.  In the random-oracle model,
if \emph{any one} of the freshly mixed per-session secrets is replaced
by an independent uniform value unknown to the adversary, then all
subsequent chaining keys and the session keys become
information-theoretically indistinguishable from uniform.

The responder-side KDF chain (Algorithm~\ref{alg:resp}) absorbs four
such secrets: $z_4 = \ECDH(e', E)$ (the ephemeral--ephemeral SM2 DH),
$\textsf{shk}_2$ (the ephemeral KEM shared secret), $z_3 =
\ECDH(e', S^i)$ (the ephemeral--static SM2 DH), and
$\textsf{shk}_3$ (the static KEM shared secret).  We give three
independent reductions, each starting from the real protocol and
replacing exactly one secret with uniform on the tested session.
Replacing $z_4$ reduces to the GDH assumption on SM2 (with $z_3$
providing redundant coverage under the same assumption); replacing
$\textsf{shk}_2$ reduces to the IND-CCA security of $\KEMe$; and
replacing $\textsf{shk}_3$ reduces to the IND-CCA security of $\KEMs$.
Since the session key becomes indistinguishable from uniform after
any one replacement, the distinguishing advantage is bounded by
$q_s$ times the minimum of the three hard-problem advantages.
The reduction sketches are given in Appendix~\ref{sec:proof}; we note
that a fully rigorous game-hopping proof---accounting for the iterative
HMAC-based KDF structure, the initiator-side KDF chain, and fine-grained
key-compromise scenarios---is omitted for brevity, but the core
arguments are standard and well-understood in the hybrid-AKE
literature~\cite{14,32}.
\end{proof}

Crucially, the theorem is \emph{parameterised by the choice of KEM}:
the reduction holds for any IND-CCA-secure KEM plugged into either the
static or ephemeral slot.  This is the computational foundation of our
agility framework---switching KEMs does not alter the security
argument, as long as the chosen KEM satisfies the IND-CCA assumption.

\subsection{Symbolic Verification with SAPIC\textsuperscript{+},
            Tamarin and ProVerif}
\label{sec:symbolic}

To complement the computational reduction, we provide a mechanised
symbolic verification of SM-PQ-WireGuard---in the tradition of
mechanised WireGuard analyses~\cite{35}---using the
SAPIC\textsuperscript{+} framework~\cite{36}.
SAPIC\textsuperscript{+} serves as a unified front-end that generates
models for multiple verification back-ends---ProVerif~\cite{37}
for reachability and equivalence properties,
Tamarin~\cite{38} for trace properties with a more
complete Diffie--Hellman equational theory, and
DeepSec~\cite{39} for confirming observational
equivalence attacks in the bounded fragment.  This multi-back-end
approach provides complementary guarantees: ProVerif is sound (but not
complete) for equivalence, Tamarin handles the full DH theory, and
DeepSec is complete for bounded equivalence.

The verification serves two purposes.  First, it confirms that the
protocol structure---the interleaving of DH, KEMs, and key
derivation---is sound even when the adversary can reveal arbitrary
combinations of long-term and ephemeral secrets.  Second, it verifies
that the agile ephemeral-KEM switching framework
(Section~\ref{sec:agility}) introduces no structural weakness, and that
every supported ephemeral KEM preserves the same security guarantees.

\subsubsection{SAPIC\textsuperscript{+} Model and Multi-Back-End
               Workflow}
\label{sec:sapicmodel}

We model the protocol in SAPIC\textsuperscript{+}'s input language,
which extends the applied $\pi$-calculus with support for
Diffie--Hellman equational theories and stateful processes.  The
classical layer (SM2 ECDH, SM3 hash, SM4-GCM/SM4-CTR) and the
post-quantum layer (KEM KeyGen, Encap, Decap) are represented as
abstract function symbols with their equational theories.  The
adversary is the standard Dolev--Yao attacker with full network
control, augmented with the ability to reveal any combination of
static keys, ephemeral keys, and encapsulation
randomness---following the minimal-compromise methodology of
Lafourcade et al.~\cite{40}.

Crucially, the model treats the SM2 group, the SM3 hash, and the KEM as
abstract symbols equipped with their equational theories, so neither
the GM/T instantiation nor the concrete KEM choice changes the
specification; the public suite identifier
(Section~\ref{sec:agility}) is a public constant available to the
adversary.  A single specification therefore covers every supported
suite, and the security properties it establishes hold for all of them.

From this single specification, SAPIC\textsuperscript{+} automatically
generates:
\begin{itemize}
\item A \textbf{ProVerif} model, used as the primary workhorse for
computing minimal compromise conditions (via its built-in support for
reachability and observational equivalence) and for the
necessary/sufficient conditions phase of the analysis.
\item A \textbf{Tamarin} model, which provides a more complete
equational theory for Diffie--Hellman and serves as an independent
cross-check of the trace properties computed by ProVerif.
\item A \textbf{DeepSec} model.  DeepSec is complete for bounded
observational equivalence, but it does not support the Diffie--Hellman
equational theory and therefore applies to the post-quantum-only
reference protocols rather than to the DH-bearing SM-PQ hybrid; for the
hybrid, ProVerif is the equivalence verifier and Tamarin independently
cross-checks the trace properties.
\end{itemize}

\subsubsection{Verification Results}
\label{sec:veriresults}

Table~\ref{tab:cnf} summarises the minimal compromise conditions
computed by ProVerif (and cross-checked with Tamarin) for the key
security properties.  Each formula is a conjunctive normal form (CNF)
over the atomic secrets (static SM2 keys, static KEM keys, ephemeral
SM2 keys, ephemeral KEM keys, and encapsulation randomness); a
property holds unless the adversary compromises a set of secrets
satisfying the CNF.

\begin{table}[htbp]
\centering
\caption{Minimal compromise conditions for key security properties.
Results are identical across ProVerif and Tamarin.  Notation:
$s^c_i, s^c_r$ = initiator/responder static SM2 keys;
$e^c_i, e^c_r$ = ephemeral SM2 keys;
$s^q_i, s^q_r$ = static KEM keys;
$e^q_i$ = ephemeral KEM key;
$r_i, r_r, r_e$ = encapsulation randomness;
$\mathsf{dh}_{s_i,s_r}$ = static-static DH secret.}
\label{tab:cnf}
\footnotesize
\tabcolsep 5pt
\begin{tabular}{ll}
\toprule
\textbf{Property} & \textbf{Minimal compromise condition (CNF)} \\
\midrule
Session uniqueness & $\varnothing$ (unconditional) \\
UKS resistance (initiator / responder) & $\varnothing$ (unconditional) \\[2pt]
Key secrecy (initiator)
  & $\mathsf{psk}\land(s^c_r\lor e^c_i)\land(\mathsf{dh}_{s_i,s_r}\lor s^c_i\lor s^c_r)\land(s^q_r\lor r_i)$ \\[2pt]
Key secrecy (responder)
  & $\mathsf{psk}\land(s^c_i\lor e^c_r)\land(\mathsf{dh}_{s_i,s_r}\lor s^c_i\lor s^c_r)\land(s^q_i\lor r_r)$ \\[2pt]
Forward secrecy
  & $\mathsf{psk}\land(s^c_r\lor e^c_i)\land(\mathsf{dh}_{s_i,s_r}\lor s^c_i\lor s^c_r)\land(e^q_i\lor r_e)\land(e^c_i\lor e^c_r)$ \\
\bottomrule
\end{tabular}
\end{table}

The key observation is that every session-key formula factors into a
classical clause and a post-quantum clause.  For key secrecy from the
initiator's perspective:
\[
\underbrace{\mathsf{psk}\land(s^c_r\lor e^c_i)\land(\mathsf{dh}_{s_i,s_r}\lor s^c_i\lor s^c_r)}_{\text{classical (SM2) clause}}
\ \land\
\underbrace{(s^q_r\lor r_i)}_{\text{post-quantum clause}}.
\]
To learn the session key, the adversary must defeat \emph{both} the
classical leg \emph{and} the post-quantum leg.  This is the precise,
machine-checked statement of hybrid security: even a quantum computer
that recovers all SM2 secrets still faces the unbroken
$(s^q_r\lor r_i)$ clause.

Crucially, these CNF formulas are \emph{independent of the concrete
KEM}.  Because the KEM enters the model only through its abstract
interface, a single ProVerif/Tamarin development establishes the
conditions for every supported suite at once: replacing the ephemeral
ML-KEM-512 with ML-KEM-768, HQC-128, or any other registered ephemeral
KEM leaves the specification---and hence the minimal
compromise conditions---unchanged.  The public suite identifier ensures
that different KEM choices produce distinct transcripts, preventing
cross-algorithm confusion, while the abstract KEM interface ensures that
the protocol structure remains unchanged.  This mechanised result
complements
Theorem~\ref{thm:hybrid}: the computational reduction guarantees
security for any IND-CCA KEM, and the symbolic verification confirms
that the protocol structure securely accommodates all of them.

The trace properties were independently re-established by Tamarin,
whose more complete Diffie--Hellman theory verified all eleven trace
lemmas---message agreement for \initmsg, \respmsg, and the confirmation
message; initiator, responder, and mutual key secrecy, each with and
without forward secrecy; and initiator and responder session
uniqueness---in agreement with the ProVerif results of
Table~\ref{tab:cnf}.

For equivalence properties (strong secrecy and anonymity), ProVerif's
observational-equivalence engine is the verifier, since DeepSec does
not support the SM2 Diffie--Hellman theory required by the hybrid.  The
headline cases---where the adversary compromises all classical secrets
(``all-dh''), all post-quantum secrets (``all-pq''), or nothing
(``no-reveal'')---return \textit{true} (observational equivalence
holds) for both anonymity and strong secrecy, confirming hybrid
security at the indistinguishability level.  In the remaining
partial-compromise scenarios ProVerif returns \textit{cannot be
proved}: its diff-equivalence over-approximation can neither certify
equivalence nor exhibit a distinguisher.  Across all evaluated
scenarios \emph{no} equivalence attack was found---of the equivalence
queries ProVerif decided, $15$ established equivalence and $19$ were
inconclusive, with zero distinguishing attacks.  We therefore claim
anonymity and strong secrecy in the no- and low-compromise regimes
where they are machine-certified, and report the higher-compromise
scenarios as not contradicted rather than proven.

\subsubsection{Properties Not Covered by Symbolic Verification}
\label{sec:notcovered}

The symbolic model does not capture side-channel resistance,
denial-of-service resistance, or the strength of the underlying
cryptographic primitives.  These are addressed as follows:
(i)~side-channel resistance is an implementation property; our
GFNI+AVX2 SM4 back-end provides a constant-time S-box on supported
platforms (Section~\ref{sec:perf}); (ii)~DoS resistance is inherited
from WireGuard's cookie-reply mechanism, with the nonce-collision
probability bounded by $2^{-53}$ (Section~\ref{sec:noncedomain}); and
(iii)~the primitive strengths are covered by the computational
assumptions in Theorem~\ref{thm:hybrid}.

\section{Performance Evaluation}
\label{sec:perf}

We evaluate SM-PQ-WireGuard on a commodity x86\_64 host to answer three
questions: (i)~what is the handshake latency and wire size overhead of
the GM/T substitution? (ii)~does the optimised SM4-GCM back-end close
the data-plane throughput gap? and (iii)~what is the runtime cost of
the agility framework itself?

\subsection{Experimental Setup}
\label{sec:setup}

All experiments run on a single core of an Intel i7-13700 (16\,GiB RAM,
Linux 6.6) compiled with Rust 1.95 in release mode with link-time
optimisation.  The GM/T primitives use standard open-source
implementations (SM2 v0.13~\cite{41}, SM3 v0.3); SM4 is
evaluated in two back-ends: the portable reference implementation (v4)
and our GFNI+AVX2 in-tree back-end (v5)~\cite{42}.  Post-quantum KEMs are provided
through the liboqs library (v0.10).  Handshake latencies are means
over $100$ handshakes and data-plane figures over $15$ micro-benchmark
rounds; the per-suite comparison of Figure~\ref{fig:kemcompare} uses
$300$ rounds with the suites interleaved round-robin so that CPU
frequency drift averages across suites rather than biasing any single
one.  Standard deviations are reported where applicable.

\subsection{Handshake Latency and Wire Size}
\label{sec:q1q2}

Table~\ref{tab:handshake} reports the handshake latency and message
sizes for three variants: stock WireGuard, the PQ-only hybrid
(PQ-WireGuard$^\star$), and SM-PQ-WireGuard.  The initiator column
measures the time to build \initmsg; the responder column measures the
time to consume \initmsg\ and build \respmsg.

\begin{table}[htbp]
\centering
\caption{Handshake latency (ms, mean over $100$ handshakes, std.\ in
parentheses) and on-wire handshake message sizes (bytes, excluding the
IPv6/UDP header).  Long-term static-KEM key pairs are generated once at
configuration time; the latencies measure the steady-state
per-handshake cost.}
\label{tab:handshake}
\footnotesize
\tabcolsep 5pt
\begin{tabular}{lrrrr}
\toprule
\textbf{Variant} & \textbf{Initiator (ms)} & \textbf{Responder (ms)}
                 & \textbf{\initmsg\ (B)} & \textbf{\respmsg\ (B)} \\
\midrule
WireGuard           & $0.083\,(0.04)$ & $0.182\,(0.08)$ & $196$  & $140$  \\
PQ-WireGuard$^\star$ & $0.144\,(0.06)$ & $0.340\,(0.13)$ & $1124$ & $1032$ \\
SM-PQ-WireGuard     & $0.621\,(0.23)$ & $1.169\,(0.40)$ & $1110$ & $1018$ \\
\bottomrule
\end{tabular}
\end{table}

Relative to an otherwise identical hybrid construction instantiated
with X25519, the only wire-format change SM-PQ-WireGuard introduces is
a single byte---the $33$-byte SEC1-compressed SM2 ephemeral public key
replacing the $32$-byte X25519 $u$-coordinate, alongside the one-byte
\texttt{f\_suite} selector.  Its messages ($1110$/$1018$\,B) are
comparable in size to PQ-WireGuard$^\star$ and, for the default
ML-KEM-512 (or Kyber-512) suite, remain within the
IPv6 MTU of $1280$ bytes (the larger suites exceed it; see below).  Every variant completes its handshake in
roughly a millisecond: with the long-term Classic~McEliece key pairs
generated once at configuration time, the per-handshake KEM work
(ML-KEM-512 key generation/encapsulation plus McEliece
encapsulation/decapsulation) accounts for only the $0.340$\,ms
PQ-WireGuard$^\star$ responder baseline.  The dominant single
contributor in the SM variant is, in fact, the \emph{classical} SM2
ECDH: it raises the responder latency to $1.169$\,ms, a
$+0.83$\,ms delta over PQ-WireGuard$^\star$ that is almost entirely
attributable to the pure-Rust SM2 implementation being markedly slower
than X25519 ($\approx 1.0$\,ms vs.\ $\approx 0.22$\,ms of ECDH time);
SM3 hashing and SM4 symmetric substitutions together add $<0.05$\,ms.
In \emph{relative} terms the GM/T substitution is therefore not free
($\approx 3\times$ the PQ-only responder time), but in \emph{absolute}
terms the entire handshake remains sub-$1.2$\,ms---imperceptible for
VPN session establishment, and dominated by a classical primitive that
a hardware SM2 engine would close.  These sub-millisecond figures
correct an earlier measurement that conflated one-time McEliece key
generation with per-handshake cost.

Because the ephemeral KEM is selected at run time through the agility
framework (Section~\ref{sec:agility}), we further measured the
per-handshake cost of \emph{every} registered suite while holding the
static McEliece anchor and the SM4-GCM data plane fixed---the KEM is the
only variable.  Figure~\ref{fig:kemcompare} reports, for the initiator
and responder respectively, the on-wire handshake size (bars, left
axis) against the handshake latency (line, right axis, $\pm$SEM over
$300$ round-robin rounds; suites are interleaved so that CPU
frequency drift averages out rather than separating the suites).

\begin{figure*}[htbp]
\centering
\begin{minipage}{0.49\linewidth}
  \centering
  \includegraphics[width=\linewidth]{kem_compare_init.png}\\
  {\footnotesize (a) Initiator (\initmsg)}
\end{minipage}\hfill
\begin{minipage}{0.49\linewidth}
  \centering
  \includegraphics[width=\linewidth]{kem_compare_respond.png}\\
  {\footnotesize (b) Responder (\respmsg)}
\end{minipage}
\caption{Per-suite cost of the agile ephemeral KEM (static McEliece
anchor and SM4-GCM data plane held fixed).  Bars: on-wire handshake
size (left axis); line: handshake latency with $\pm$SEM error bars
(right axis, $300$ rounds).  On-wire size spans nearly $10\times$
(ML-KEM-512 $\approx 1.0$\,KiB to FrodoKEM-640 $\approx 9.7$\,KiB) and
is deterministic, whereas latency separates only the code-based
HQC-128; the remaining five suites (including the $9.7$\,KiB FrodoKEM)
are statistically indistinguishable in time.}
\label{fig:kemcompare}
\end{figure*}

Two observations follow.  First, size and speed are
\emph{uncorrelated} across the agile suites: HQC-128 is the latency
outlier ($3.18$\,ms responder) at a moderate $4.6$\,KiB \respmsg, while
FrodoKEM-640 is the size outlier ($9.7$\,KiB, nearly $10\times$
ML-KEM-512) yet sits in the same latency cluster as the lattice suites
($1.35$\,ms responder)---a huge ciphertext at ordinary speed.  Second,
the five non-code suites (ML-KEM-512/768/1024, Kyber-512, and
FrodoKEM-640) span only $1.12$--$1.35$\,ms at the responder, a range
narrower than their own run-to-run standard deviation, and are thus
distinguishable only by wire size---which makes message size, not time,
the operative cost metric when an operator chooses a security level.
This is precisely the trade-off space the agility framework exposes
through its single external selector: an operator can move along the
NIST-level axis (L1$\to$L5) or hedge across algorithm families
(lattice/code/LWE) at deployment time, paying in bytes and, for the
code-based option, in roughly $2$\,ms of extra latency.  Concretely,
only the default ML-KEM-512 and Kyber-512 suites keep both handshake
messages within the IPv6 MTU of $1280$\,bytes (G4); the larger suites
(ML-KEM-768/1024 at $\approx 1.3$--$1.9$\,KiB, HQC-128 at up to
$4.6$\,KiB, and Frodo-640-AES at $\approx 9.7$\,KiB) exceed it and are
delivered via IP fragmentation, which is the explicit cost an operator
accepts when selecting a higher security level or a non-lattice family.

\subsection{Data-Plane Throughput}
\label{sec:q5}

The data plane is the one place where the GM/T substitution carries a
real runtime cost.  Table~\ref{tab:throughput} reports single-core
encryption throughput for four implementations across
WireGuard-relevant packet sizes: the optimised assembly
ChaCha20--Poly1305 implementation (the stock WireGuard back-end),
the portable software ChaCha20--Poly1305 implementation (software baseline),
portable SM4-GCM (v4), and our GFNI+AVX2 SM4-GCM (v5).

\begin{table}[htbp]
\centering
\caption{Single-core AEAD encryption throughput (MiB/s).  v5 is our
GFNI+AVX2 SM4-GCM back-end.}
\label{tab:throughput}
\footnotesize
\tabcolsep 5pt
\begin{tabular}{rrrrrr}
\toprule
\textbf{Payload (B)} & \textbf{ring (ChaCha20)} & \textbf{soft (ChaCha20)}  & \textbf{SM4-GCM v4} & \textbf{SM4-GCM v5} & \textbf{v5/v4} \\
\midrule
$64$    & $417$  & $59$  & $56$  & $98$  & $1.7\times$ \\
$512$   & $1592$ & $302$ & $100$ & $383$ & $3.8\times$ \\
$1280$  & $2230$ & $432$ & $105$ & $517$ & $4.9\times$ \\
$1500$  & $2319$ & $481$ & $97$  & $497$ & $5.2\times$ \\
$8192$  & $2834$ & $618$ & $109$ & $605$ & $5.6\times$ \\
$65536$ & $2642$ & $715$ & $99$  & $618$ & $6.3\times$ \\
\bottomrule
\end{tabular}
\end{table}

Our GFNI+AVX2 SM4-GCM back-end (v5) is $4.4$--$6.3\times$ faster than
the portable v4, with the speed-up growing with packet size as the
per-call fixed cost (key schedule, the CLMUL-based GHASH
subkey~\cite{43}) is amortised and SIMD lanes fill.  At WireGuard MTUs ($1280$--$1500$\,B), v5 reaches parity
with the portable software ChaCha20--Poly1305 implementation
($0.98$--$1.20\times$).  The residual
$\sim 4.5\times$ gap against \texttt{ring}'s hand-written AVX2 assembly
is the only remaining headroom, addressable by a hardware SM4 engine or
future hand-tuned assembly.

\subsection{Agility Framework Overhead}
\label{sec:agilityoverhead}

To quantify the cost of cryptographic agility itself, we compare three
configurations of the same binary: (i)~no agility framework (hard-coded
KEM selection); (ii)~agility framework with KEM selection resolved at
compile time; and (iii)~agility framework with KEM selection resolved
at deployment time via configuration.  All three use identical
cryptographic primitives (SM2 + ML-KEM-512).

\begin{table}[htbp]
\centering
\caption{Agility framework overhead (responder latency, ms; ML-KEM-512,
$100$ handshakes).  Only the deployment-time row is separately
benchmarked (the measured agile binary); the upper two rows ($\approx$)
are equal to it by construction and not independently measured, since
the dispatch cost (an atomic load plus a \texttt{match}) lies below the
per-handshake standard deviation.}
\label{tab:agility}
\footnotesize
\tabcolsep 10pt
\begin{tabular}{lr}
\toprule
\textbf{Configuration} & \textbf{Responder latency (ms)} \\
\midrule
No agility (hard-coded)     & $\approx 1.169$ \\
Compile-time dispatch       & $\approx 1.169$ \\
Deployment-time dispatch    & $1.169\,(0.40)$ \\
\bottomrule
\end{tabular}
\end{table}

The agility framework's only cost on the wire is the single
\texttt{f\_suite} byte; at run time, suite selection is an atomic load
and a \texttt{match}, and binding the \texttt{f\_suite} byte into the
transcript adds one SM3 compression---together a few nanoseconds, far
below the
$\pm0.40$\,ms run-to-run standard deviation of the handshake itself.
The framework therefore imposes no measurable handshake overhead: the
flexibility to switch among ML-KEM-512/768/1024, Kyber-512, HQC-128 and
FrodoKEM-640 (Figure~\ref{fig:kemcompare}) is effectively free relative
to the cryptographic work it dispatches.

\subsection{Summary}
\label{sec:perfsummary}

In summary, SM-PQ-WireGuard adds \textbf{one byte} to the wire format
(relative to an X25519 hybrid) and keeps the entire handshake
\textbf{sub-$1.2$\,ms}.  The GM/T
substitution's cost is concentrated in the classical SM2 ECDH
($\approx 0.83$\,ms over the PQ-only hybrid)---a large relative but
negligible absolute overhead, and one orthogonal to the post-quantum
KEMs, whose per-handshake cost is only $\approx 0.34$\,ms.  The
data-plane gap is closed by the GFNI+AVX2 v5 back-end, reaching
portable software ChaCha20-Poly1305 parity at WireGuard MTUs.  The
agility framework itself adds no measurable handshake overhead beyond
the single wire byte, validating that flexible KEM switching---across
ML-KEM-512/768/1024, Kyber-512, HQC-128, and FrodoKEM-640-AES, whose
per-suite costs differ by up to $10\times$ in wire size while the
handshake stays near $1$\,ms for every suite except the code-based
HQC-128 ($\approx 3.2$\,ms) (Figure~\ref{fig:kemcompare})---comes
essentially for free.

% =====================================================================
\section{Discussion and Conclusion}
\label{sec:disc}

\subsection{Discussion}

Our results show that SM-PQ-WireGuard provides a practical path toward
a regulation-compliant, post-quantum-secure VPN handshake with minimal
performance overhead.  Several aspects merit further discussion.

\paragraph{Hardware acceleration.}
With the SM4 data-plane gap closed by our GFNI+AVX2 back-end, the
single largest remaining GM/T cost is the SM2 scalar multiplication,
which accounts for approximately $1$\,ms of cryptographic time per
handshake (Table~\ref{tab:handshake}) and is $3.8\times$ slower than
X25519 in the atomic benchmark.  This gap is not intrinsic to SM2; it
reflects the absence of optimised, vectorised scalar multiplication in
the current SM2 software implementation.  Many newer Chinese SoCs implement
SM2 (and SM4) in hardware; on such platforms, we expect the
SM-PQ-WireGuard handshake to fall back toward the
PQ-WireGuard$^\star$ baseline, with the SM2 contribution reduced to
near-X25519 levels and the GM/T overhead below $1\%$.

\paragraph{Side-channel considerations.}
The v5 data-plane SM4 provides a constant-time S-box on GFNI-capable
hosts by computing the S-box via the GFNI affine instruction pair,
which perform a fixed Galois-field affine map with no table look-ups
and hence no data-dependent memory access---the classic source of
cache- and timing-based key recovery~\cite{44,45}.  The residual side-channel
risk is in the SM2 scalar multiplication: the windowed implementation
in the SM2 software library is not as battle-tested as X25519's
Montgomery ladder.  A constant-time SM2 back-end---either via a
bit-sliced implementation or a hardware engine---is the most important
hardening priority for GM/T-compliant deployments.

\paragraph{KEM selection and future-proofing.}
The agility framework supports a range of KEM options
(ML-KEM-512/768/1024, Kyber-512, HQC-128, Frodo-640-AES).  Our
performance evaluation (Section~\ref{sec:perf}) shows that the runtime
difference between KEMs is dominated by the KEM operations themselves,
not by the agility framework, so operators can choose based on security
margins and key-size constraints rather than performance.  The
ProVerif-verified domain separation guarantees that switching KEMs does
not introduce cross-algorithm confusion or invalidate the security
proof.  As the NIST post-quantum standardisation landscape evolves, new
KEMs can be added to the framework by implementing the same abstract
interface; no protocol-level changes are required.

\paragraph{Other future work.}
Additional directions include refining the cookie-reply nonce field
(reclaiming $12$ bytes per cookie reply), integrating a hardware SM2
back-end, and extending the constant-time SM4 path to non-GFNI
platforms via a bit-sliced fallback or a dedicated ISA
extension~\cite{46}.

\subsection{Conclusion}

We presented SM-PQ-WireGuard, a hybrid post-quantum WireGuard handshake
that unifies GM/T compliance with cryptographic agility.  The protocol
instantiates the classical layer with SM2, SM3, and SM4, and
combines a fixed Classic-McEliece-460896 static-KEM identity anchor
with an agile ephemeral-KEM slot supporting a range of post-quantum
options including ML-KEM, Kyber, HQC, and FrodoKEM.  An
algorithm-agnostic abstraction
layer and a one-byte suite-identifier domain separator enable flexible
switching of the ephemeral KEM without protocol-level changes or
security-proof invalidation.  We explicitly re-derived the Noise transcript for the
GM/T setting, resolving two protocol-layer mismatches---MAC1/MAC2
reconstruction via HMAC-SM3 truncation and per-site nonce partitioning
for SM4-GCM---and provided a dual-track security analysis combining an
IND-CCA reduction with mechanised ProVerif/Tamarin verification.  Our
implementation adds one byte to the wire format relative to an X25519
hybrid, keeps the whole handshake sub-$1.2$\,ms (the GM/T overhead being
concentrated in the
classical SM2 ECDH rather than the post-quantum KEMs), and reaches
portable software ChaCha20-Poly1305 parity at WireGuard MTUs with a GFNI+AVX2 SM4
back-end; the agility framework itself contributes no measurable
handshake overhead beyond the single suite-selector byte.  We believe SM-PQ-WireGuard is a credible deployment target
for the Chinese post-quantum VPN ecosystem, offering a practical
balance of GM/T standards compliance, quantum security, and cryptographic
flexibility.

% =====================================================================
%%% Acknowledgements
%%%%%%%%%%%%%%%%%%%%%%%%%%%%%%%%%%%%%%%%%%%%%%%%%%%%%%%
\Acknowledgements{This work was supported by Zhejiang Leading Innovation and Entrepreneurship Team (No.~2024R02004).}

% =====================================================================
%%% Supplements
%%%%%%%%%%%%%%%%%%%%%%%%%%%%%%%%%%%%%%%%%%%%%%%%%%%%%%%
\Supplements{Appendix~A contains a proof sketch of
Theorem~\ref{thm:hybrid}.  The supporting information is available
online at info.scichina.com and link.springer.com.  The supporting
materials are published as submitted, without typesetting or editing.
The responsibility for scientific accuracy and content remains entirely
with the authors.}

% =====================================================================
%%% References
%%%%%%%%%%%%%%%%%%%%%%%%%%%%%%%%%%%%%%%%%%%%%%%%%%%%%%%
\begin{thebibliography}{99}

\bibitem{1} Donenfeld J A. WireGuard: Next Generation Kernel Network Tunnel. In: Proceedings of NDSS, San Diego, 2017

\bibitem{2} Perrin T. The Noise Protocol Framework, rev. 34. 2018. \url{https://noiseprotocol.org/noise.html}

\bibitem{3} Diffie W, Hellman M E. New directions in cryptography. IEEE Trans Inform Theory, 1976, 22: 644--654

\bibitem{4} Bernstein D J. Curve25519: new Diffie--Hellman speed records. In: Proceedings of PKC, New York, 2006. 207--228

\bibitem{5} Nir Y, Langley A. ChaCha20 and Poly1305 for IETF protocols. IETF RFC 8439, 2018

\bibitem{6} Aumasson J-P, Neves S, Wilcox-O'Hearn Z, et al. BLAKE2: simpler, smaller, fast as MD5. In: Proceedings of ACNS, Banff, 2013. 119--135

\bibitem{7} Miller V S. Use of elliptic curves in cryptography. In: Proceedings of CRYPTO, Santa Barbara, 1985. 417--426

\bibitem{8} Koblitz N. Elliptic curve cryptosystems. Math Comput, 1987, 48: 203--209

\bibitem{9} Shor P W. Algorithms for quantum computation: discrete logarithms and factoring. In: Proceedings of FOCS, Santa Fe, 1994. 124--134

\bibitem{10} Mosca M. Cybersecurity in an era with quantum computers: will we be ready? IEEE Secur Priv, 2018, 16: 38--41

\bibitem{11} Stebila D, Fluhrer S, Gueron S. Hybrid key exchange in TLS~1.3. IETF Internet-Draft, 2020

\bibitem{12} Bos J W, Costello C, Naehrig M, et al. Post-quantum key exchange for the TLS protocol from the ring learning with errors problem. In: Proceedings of IEEE S\&P, San Jose, 2015. 553--570

\bibitem{13} Bindel N, et al. Hybrid key encapsulation mechanisms and authenticated key exchange. In: Proceedings of PQCrypto, Chongqing, 2019. 206--225

\bibitem{14} Dowling B, Paterson K G. A cryptographic analysis of the WireGuard protocol. In: Proceedings of ACNS, Leuven, 2018. 3--21

\bibitem{15} NIST. FIPS 203: Module-Lattice-Based Key-Encapsulation Mechanism Standard. 2024

\bibitem{16} Albrecht M R, et al. Classic McEliece: conservative code-based cryptography. NIST PQC Standardization, 2022

\bibitem{17} State Cryptography Administration. GM/T 0003.1-2012: SM2, Part~1: General. 2012

\bibitem{18} State Cryptography Administration. GM/T 0003.5-2012: SM2, Part~5: Parameter definition. 2012

\bibitem{19} State Cryptography Administration. GM/T 0004-2012: SM3 cryptographic hash algorithm. 2012

\bibitem{20} SAC, AQSIQ. GB/T 32907-2016: SM4 block cipher algorithm. 2016

\bibitem{21} Katz J, Lindell Y. Introduction to Modern Cryptography. 2nd ed. CRC Press, 2014

\bibitem{22} Rogaway P. Authenticated-encryption with associated-data. In: Proceedings of CCS, Washington, 2002. 98--107

\bibitem{23} Bellare M, Namprempre C. Authenticated encryption: relations among notions. In: Proceedings of ASIACRYPT, Kyoto, 2000. 531--545

\bibitem{24} Certicom Research. SEC 1: Elliptic Curve Cryptography, v2.0. 2009

\bibitem{25} Merkle R C. A certified digital signature. In: Proceedings of CRYPTO, Santa Barbara, 1989. 218--238

\bibitem{26} Damg\aa rd I B. A design principle for hash functions. In: Proceedings of CRYPTO, Santa Barbara, 1989. 416--427

\bibitem{27} Dworkin M. NIST SP 800-38D: Recommendation for Block Cipher Modes of Operation: Galois/Counter Mode (GCM) and GMAC. 2007

\bibitem{28} McGrew D A, Viega J. The security and performance of the Galois/Counter Mode (GCM). In: Proceedings of INDOCRYPT, Chennai, 2004. 343--355

\bibitem{29} Avanzi R, et al. CRYSTALS-Kyber (v3.02). NIST PQC Standardization, 2022

\bibitem{30} Krawczyk H, Bellare M, Canetti R. HMAC: keyed-hashing for message authentication. IETF RFC 2104, 1997

\bibitem{31} Krawczyk H. SIGMA: the ``SIGn-and-MAc'' approach to authenticated Diffie--Hellman and its use in the IKE protocols. In: Proceedings of CRYPTO, Santa Barbara, 2003. 399--424

\bibitem{32} Shoup V. Sequences of games: a tool for taming complexity in security proofs. IACR Cryptology ePrint Archive, 2004, 2004/332

\bibitem{33} Bellare M, Rogaway P. Random oracles are practical. In: Proceedings of CCS, Fairfax, 1993. 62--73

\bibitem{34} Bellare M. New proofs for NMAC and HMAC: security without collision-resistance. In: Proceedings of CRYPTO, Santa Barbara, 2006. 602--619

\bibitem{35} Lipp B, Blanchet B, Bhargavan K. A mechanised cryptographic proof of the WireGuard VPN protocol. In: Proceedings of IEEE EuroS\&P, Stockholm, 2019. 231--246

\bibitem{36} Cheval V, Jacomme C, Kremer S, et al. SAPIC\textsuperscript{+}: protocol verifiers of the world, unite! In: Proceedings of 31st USENIX Security Symposium, Boston, 2022. 3935--3952

\bibitem{37} Blanchet B. Modeling and verifying security protocols with the applied pi calculus and ProVerif. Found Trends Priv Secur, 2016, 1: 1--135

\bibitem{38} Meier S, Schmidt B, Cremers C, et al. The TAMARIN prover for the symbolic analysis of security protocols. In: Proceedings of CAV, Saint Petersburg, 2013. 696--701

\bibitem{39} Cheval V, Kremer S, Rakotonirina I. DEEPSEC: deciding equivalence properties in security protocols. In: Proceedings of IEEE S\&P, San Francisco, 2018. 542--560

\bibitem{40} Lafourcade P, Mahmoud D, Ruhault S, et al. A tale of two worlds: a formal story of WireGuard hybridization. In: Proceedings of 34th USENIX Security Symposium, Seattle, 2025. 4937--4956

\bibitem{41} RustCrypto Project. sm2 crate, version~0.13. \url{https://crates.io/crates/sm2}

\bibitem{42} Guo W. Efficient constant-time implementation of SM4 with Intel GFNI instruction set extension and Arm NEON coprocessor. IACR Cryptology ePrint Archive, 2022, 2022/1154

\bibitem{43} Gueron S, Kounavis M E. Intel carry-less multiplication instruction and its usage for computing the GCM mode. Intel White Paper, 2010 (rev. 2014)

\bibitem{44} Kocher P C. Timing attacks on implementations of Diffie--Hellman, RSA, DSS, and other systems. In: Proceedings of CRYPTO, Santa Barbara, 1996. 104--113

\bibitem{45} Bernstein D J. Cache-timing attacks on AES. Technical report, 2005

\bibitem{46} Saarinen M-J O. A lightweight ISA extension for AES and SM4. arXiv:2002.07041, 2020

\end{thebibliography}

% =====================================================================
%%% Appendix
%%%%%%%%%%%%%%%%%%%%%%%%%%%%%%%%%%%%%%%%%%%%%%%%%%%%%%%
\begin{appendix}

\section{Proof Sketch of Theorem~\ref{thm:hybrid}}
\label{sec:proof}

We give the core reductions in detail while explicitly acknowledging
where formal game hops are elided for brevity.  A fully rigorous proof
would expand each step into several sub-games (accounting for the
iterative HMAC-based KDF, bad-flag analyses for ROM consistency, and
fine-grained key-compromise cases) and would additionally reduce the
initiator-side KDF chain ($z_1$, $\textsf{shk}_1$, $z_2$)---all of
which is standard in the hybrid-AKE literature~\cite{14,32,35} and
omitted here.

\subsection{ROM Setup and KDF Chain}

We model SM3 as a random oracle $\mathcal{O} \colon \{0,1\}^* \to
\{0,1\}^{256}$.  HMAC-SM3 is treated as a keyed RO via the standard
NMAC reduction~\cite{34}.  The KDF of Section~\ref{sec:handshake}
is a composition of RO calls; for brevity we write each $\KDF_n(ck,m)$
as a sequence of $\mathcal{O}$ queries and focus on the information-flow
structure rather than the exact query count.

After both parties have processed \initmsg\ they share a chaining key
$ck_{\text{in}}$.  The remaining KDF chain (common to both parties,
from Algorithm~\ref{alg:resp}) is:
\begin{align}
ck_1 &\gets \KDF_1(ck_{\text{in}},\; \pkenc(E') \concat \textit{ct}_2)
  \quad\text{(public input)}, \label{eq:apx1} \\
ck_2 &\gets \KDF_1(ck_1,\; z_4 \concat \textsf{shk}_2), \label{eq:apx2} \\
ck_3 &\gets \KDF_1(ck_2,\; z_3 \concat \textsf{shk}_3), \label{eq:apx3} \\
(ck_4, \tau, k_E) &\gets \KDF_3(ck_3,\; Q), \label{eq:apx4} \\
(k_{r\to i},\, k_{i\to r}) &\gets \KDF_2(ck_4,\; \varepsilon), \label{eq:apx5}
\end{align}
where $z_3 = \ECDH(e', S^i)$, $z_4 = \ECDH(e', E)$,
$\textsf{shk}_2 = \Decap(e_q, \textit{ct}_2)$, and
$\textsf{shk}_3 = \Decap(\mathit{sk}_{S^i_q}, \textit{ct}_3)$.

In the ROM, if any one of $\{z_4, \textsf{shk}_2, z_3, \textsf{shk}_3\}$
is replaced by an independent uniform value, the corresponding
$\mathcal{O}$ query receives a fresh $256$-bit suffix unknown to the
adversary.  The RO output---and all subsequent values in the
chain---are therefore indistinguishable from uniform (the adversary
can guess the fresh suffix only with probability
$q_{\mathcal{O}}/2^{256}$, which is negligible).  Consequently, the
session keys $k_{r\to i}$ and $k_{i\to r}$ are pseudorandom if
\emph{any single} one of these four secrets remains hidden.

$z_3$ and $z_4$ are both SM2 DH values; replacing either independently
suffices.  We give the reduction for $z_4$ (the ephemeral--ephemeral
DH) as the classical leg; $z_3$ is covered by an identical argument.

\subsection{Path A: Classical DH (GDH)}

Replace $z_4$ with uniform $r_4 \xleftarrow{\$} \{0,1\}^{256}$ on the
tested session.  A distinguisher between real $z_4$ and uniform $r_4$
yields a GDH solver:

Given a GDH challenge $(A=aG,\, B=bG)$ with a DDH oracle
$\mathcal{O}_{\text{DDH}}$, the reduction guesses the test session
(probability $1/q_s$) and embeds:
\begin{itemize}
\item Initiator's ephemeral SM2 public key: $E \gets A = aG$;
\item Responder's ephemeral SM2 public key: $E' \gets B = bG$.
\end{itemize}
The reduction does \emph{not} know $a$ or $b$, but using the known
static private keys it computes all other required DH values:
$z_1 = \mathit{sk}_{S^r} \!\cdot\! A$,
$z_2 = \ECDH(\mathit{sk}_{S^i}, S^r)$,
$z_3 = \mathit{sk}_{S^i} \!\cdot\! B$.
The target $z_4 = abG$ is the CDH solution.  The reduction simulates
the protocol using $r_4$ in place of $z_4$ and answers RO queries
with a lazy-sampling table.

If the adversary queries $\mathcal{O}$ on an input containing $abG$,
the reduction detects it via $\mathcal{O}_{\text{DDH}}(A,B,\cdot)$ and
outputs the CDH solution.  If no such query occurs, the real and
simulated views are identically distributed.  Hence
\[
|\Adv_{\text{real}} - \Adv_{\text{rand}-z_4}| \le
q_s \!\cdot\! \Adv^{\mathrm{GDH}}_{\textsf{SM2}}(\mathcal{B}_1)
+ \negl(\lambda).
\]
(We elide the formal game hop $\textsc{Game}_0 \to \textsc{Game}_1$
that enforces ephemeral-key collision-freeness and the hop
$\textsc{Game}_1 \to \textsc{Game}_2$ that programs the RO; both are
standard and contribute only $\negl(\lambda)$ terms.)

\subsection{Path B: Ephemeral KEM (IND-CCA)}

Replace $\textsf{shk}_2$ with uniform.  The reduction embeds an
IND-CCA challenge $(\mathit{pk}^\star, c^\star, k^\star_d)$ against
$\KEMe$: set $E_q \gets \mathit{pk}^\star$, $\textit{ct}_2 \gets
c^\star$, $\textsf{shk}_2 \gets k^\star_d$.  By KEM correctness,
$k^\star_0 = \Decap(\mathit{sk}^\star, c^\star)$, so using $k^\star_d$
in place of the decapsulation output is consistent with the protocol.
The IND-CCA decapsulation oracle handles decapsulation queries for
other sessions.  If $d=0$ (real KEM secret) the view is the real
protocol; if $d=1$ (uniform) the view matches the $\textsf{shk}_2$-replaced
game.  Thus
\[
|\Adv_{\text{real}} - \Adv_{\text{rand}-\textsf{shk}_2}| \le
q_s \!\cdot\! \Adv^{\mathrm{IND\text{-}CCA}}_{\KEMe}(\mathcal{B}_2)
+ \negl(\lambda).
\]
(A formal proof would add game hops for the KEM key-pair collision and
for ciphertext-collision events across sessions; both are bounded by
$q_s^2/2^{256}$.)

\subsection{Path C: Static KEM (IND-CCA)}

Replace $\textsf{shk}_3$ with uniform.  The structure is identical to
Path B, embedding the IND-CCA challenge for $\KEMs$: set the initiator's
static KEM public key $S^i_q \gets \mathit{pk}^\star$, and in the test
session set $\textit{ct}_3 \gets c^\star$, $\textsf{shk}_3 \gets
k^\star_d$.  The stream-cipher wrapping $\textit{cct}_3$ is simulated
honestly using the known $z_3$.  Other sessions that use $S^i_q$ are
handled via the IND-CCA decapsulation oracle.  We obtain
\[
|\Adv_{\text{real}} - \Adv_{\text{rand}-\textsf{shk}_3}| \le
q_s \!\cdot\! \Adv^{\mathrm{IND\text{-}CCA}}_{\KEMs}(\mathcal{B}_3)
+ \negl(\lambda).
\]

\subsection{Combining the Bounds}

After replacing any one of $z_4$, $\textsf{shk}_2$, or $\textsf{shk}_3$
with uniform, the session keys are indistinguishable from uniform in
the ROM (by the argument of Section~A.1 applied to
the chain~\eqref{eq:apx2}--\eqref{eq:apx5}).  Hence each path gives an
independent upper bound on the real-protocol advantage.  Taking the
tightest (minimum) of the three bounds and noting that SM4 key
truncation ($256 \to 128$ bits) preserves uniformity yields
\[
\Adv^{\textsc{ror}}_{\Pi}(\mathcal A) \le
q_s\cdot\min\bigl\{
\Adv^{\mathrm{GDH}}_{\textsf{SM2}},\;
\Adv^{\mathrm{IND\text{-}CCA}}_{\KEMe},\;
\Adv^{\mathrm{IND\text{-}CCA}}_{\KEMs}
\bigr\} + \negl(\lambda),
\]
which is Theorem~\ref{thm:hybrid}.  The minimum formalises hybrid
security: the protocol remains secure as long as \emph{any one} of the
classical SM2 leg, the ephemeral post-quantum leg, or the static
post-quantum leg resists cryptanalysis.

\subsection*{Omitted Details}

A fully rigorous version of this proof would additionally:
(i)~expand each $\KDF_n$ call into its constituent HMAC-based RO
queries with explicit domain separation, yielding
$|\Adv_i - \Adv_{i+1}|$ bounds for each intermediate game;
(ii)~analyse the initiator-side KDF chain ($z_1$, $\textsf{shk}_1$,
$z_2$ absorbed into $ck_{\text{in}}$), which provides an independent
earlier barrier and would tighten the final bound;
(iii)~formalise the hybrid freshness condition
$\textsf{F1} \lor \textsf{F2} \lor \textsf{F3}$ in the RoR experiment;
(iv)~add bad-flag analyses for RO collisions and ciphertext collisions.
All of these are routine in the game-hopping
framework~\cite{32,35} and their omission does not affect the
validity of the core reductions presented above.

\end{appendix}

\end{document}